\begin{table*}[t]
\centering
\caption{Architectural Comparison of Evaluated Linux Inter-Process Communication Mechanisms}
\label{tab:ipc_comparison}
\scriptsize
\begin{tabular}{lp{2.4cm}ccccp{3.2cm}p{3.2cm}}
\toprule
\textbf{Mechanism} & \textbf{Data Structure / Model} & \textbf{Data Copies} & \textbf{Buffering Space} & \textbf{Boundary Semantics} & \textbf{Max Payload Limit} & \textbf{Synchronization Primitive} & \textbf{Key System Calls} \\
\midrule
\textbf{Anonymous Pipe} & Kernel page ring buffer & 2 & Kernel & Byte-stream & Bounded by pipe buffer (64 KB) & Implicit (kernel blocking queue) & \texttt{pipe2}, \texttt{read}, \texttt{write}, \texttt{close} \\
\textbf{Named FIFO} & Filesystem inode pipe & 2 & Kernel & Byte-stream & Bounded by pipe buffer (64 KB) & Implicit (kernel blocking queue) & \texttt{mkfifo}, \texttt{open}, \texttt{read}, \texttt{write} \\
\textbf{System V MQ} & Kernel linked message list & 2 & Kernel & Message-oriented & \texttt{MSGMAX} / \texttt{MSGMNB} sysctl & Implicit (message type filtering) & \texttt{msgget}, \texttt{msgsnd}, \texttt{msgrcv}, \texttt{msgctl} \\
\textbf{POSIX MQ} & Virtual \texttt{/dev/mqueue} file & 2 & Kernel & Priority message & Dynamic \texttt{mq\_msgsize} (max 10 msgs) & Implicit (priority queue lock) & \texttt{mq\_open}, \texttt{mq\_send}, \texttt{mq\_receive} \\
\textbf{POSIX SHM} & Memory-mapped pages & \textbf{0} & Userspace & Structured mailbox & 1,024 B (\texttt{SHM\_MAX\_PAYLOAD}) & Explicit (POSIX unnamed semaphores) & \texttt{shm\_open}, \texttt{mmap}, \texttt{sem\_wait}, \texttt{sem\_post} \\
\textbf{Unix Domain Skt} & Socket buffer (\texttt{sk\_buff}) & 2 & Kernel & Byte-stream & Socket buffer auto-tuning & Implicit (BSD socket transport) & \texttt{socket}, \texttt{bind}, \texttt{connect}, \texttt{send}, \texttt{recv} \\
\textbf{Kernel Mailbox} & Custom char device & 2 & Kernel & Bounded byte-stream & 64 KB ring capacity & Sleepable mutex + wait queues & \texttt{open}, \texttt{read}, \texttt{write}, \texttt{poll}, \texttt{release} \\
\bottomrule
\end{tabular}
\end{table*}
